\documentclass[12pt,a4paper,openany,oneside]{book}

% --- Paquetes ---
\usepackage[utf8]{inputenc}
\usepackage[spanish, es-noshorthands]{babel}
\usepackage{amsmath, amsfonts, amssymb}
\usepackage{graphicx}
\usepackage{geometry}
\usepackage{hyperref}
\usepackage{booktabs}
\usepackage{fancyhdr}
\usepackage{listings}
\usepackage{xcolor}
\usepackage{tikz}
\usepackage{pgfplots}
\usetikzlibrary{arrows.meta, positioning, shapes.geometric}
\pgfplotsset{compat=1.18}

\geometry{margin=2.5cm}

% --- Colores y estilos para código Python ---
\definecolor{codegreen}{rgb}{0,0.6,0}
\definecolor{codegray}{rgb}{0.5,0.5,0.5}
\definecolor{codepurple}{rgb}{0.58,0,0.82}
\definecolor{backcolour}{rgb}{0.95,0.95,0.92}

\lstset{
    language=Python,
    backgroundcolor=\color{backcolour},   
    commentstyle=\color{codegreen},
    keywordstyle=\color{blue},
    numberstyle=\tiny\color{codegray},
    stringstyle=\color{codepurple},
    basicstyle=\ttfamily\small,
    breaklines=true,
    frame=single,
    showstringspaces=false
}

% --- Estilo de cabecera y pie ---
\pagestyle{fancy}
\fancyhf{}
\renewcommand{\headrulewidth}{0.4pt}
\renewcommand{\footrulewidth}{0.4pt}
\fancyhead[L]{\small Carlos César Sánchez Coronel}
\fancyhead[R]{\small Sesión 9: Automatizando el Gradiente}
\fancyfoot[L]{\small Carlos César Sánchez Coronel}
\fancyfoot[C]{\thepage}

% --- Portada ---
\title{
    \textbf{Sesión 9} \\ 
    \Huge \textbf{AUTOMATIZANDO EL GRADIENTE} \\
    \large Introducción a PyTorch/TensorFlow y Descenso de Gradiente \\
    \normalsize Diferenciación automática, grafos computacionales y optimizadores modernos
}
\author{
    \small Por \\ \vspace{0.2cm}
    \Large \textsc{Carlos César Sánchez Coronel}
}
\date{2026}

\begin{document}

\maketitle
\tableofcontents
\addcontentsline{toc}{chapter}{Índice general}

% ------------------------------------------------------------------
% CAPÍTULO 9: SESIÓN 9 - AUTOMATIZANDO EL GRADIENTE
% ------------------------------------------------------------------
\chapter{Automatizando el Gradiente: Introducción a PyTorch/TensorFlow y Descenso de Gradiente}

En las sesiones anteriores aprendimos a calcular derivadas y gradientes manualmente, y comprendimos que el descenso de gradiente es el corazón del entrenamiento de modelos de inteligencia artificial. Sin embargo, en la práctica, los modelos tienen millones de parámetros y sería imposible derivar a mano todas las expresiones. Aquí entran en juego los frameworks de deep learning como PyTorch y TensorFlow, que implementan la \textbf{diferenciación automática} (\textit{automatic differentiation}) y construyen \textbf{grafos computacionales} para calcular gradientes de manera eficiente. En esta sesión exploraremos estos conceptos, implementaremos descenso de gradiente desde cero con NumPy, y luego veremos cómo los frameworks modernos automatizan todo el proceso.

\section{Descenso de gradiente: recordatorio y variantes}

\subsection{El algoritmo básico}
Recordemos que para minimizar una función de pérdida $\mathcal{L}(\mathbf{w})$ dependiente de unos parámetros $\mathbf{w}$, el descenso de gradiente actualiza los parámetros en la dirección opuesta al gradiente:
\[
\mathbf{w}_{t+1} = \mathbf{w}_t - \eta \nabla \mathcal{L}(\mathbf{w}_t)
\]
donde $\eta$ es la \textbf{tasa de aprendizaje} (learning rate) y $t$ es el número de iteración.

\subsection{Épocas, batches y variantes}
En machine learning trabajamos con conjuntos de datos. La forma de presentar los datos al algoritmo da lugar a distintas variantes:

\begin{itemize}
    \item \textbf{Batch Gradient Descent:} Se calcula el gradiente usando \textbf{todo} el conjunto de datos. Es costoso para grandes datasets.
    \[
    \nabla \mathcal{L}(\mathbf{w}) = \frac{1}{N} \sum_{i=1}^{N} \nabla \ell_i(\mathbf{w})
    \]
    \item \textbf{Stochastic Gradient Descent (SGD):} Se usa una sola muestra aleatoria por iteración. Es rápido pero ruidoso.
    \[
    \nabla \mathcal{L}(\mathbf{w}) \approx \nabla \ell_i(\mathbf{w}) \quad \text{para un } i \text{ aleatorio}
    \]
    \item \textbf{Mini-batch Gradient Descent:} Compromiso entre ambos: se usa un subconjunto (mini-batch) de tamaño $B$.
    \[
    \nabla \mathcal{L}(\mathbf{w}) \approx \frac{1}{B} \sum_{i \in \mathcal{B}} \nabla \ell_i(\mathbf{w})
    \]
\end{itemize}

Una \textbf{época} (epoch) corresponde a haber recorrido todo el conjunto de datos una vez.

\section{Diferenciación automática y grafos computacionales}

\subsection{Comparativa de métodos de diferenciación}
Antes de profundizar en la diferenciación automática, es útil comparar las tres rutas principales para calcular derivadas en un computador:

\begin{table}[h]
\centering
\begin{tabular}{@{}llll@{}}
\toprule
\textbf{Método} & \textbf{Exactitud} & \textbf{Costo Computacional} & \textbf{Ventaja Principal} \\ \midrule
Numérico (diferencias finitas) & Baja & Medio & Fácil implementación \\
Simbólico (SymPy) & Exacta & Muy Alto & Expresión analítica \\
Automático (Autodiff) & Exacta & Bajo & Eficiencia en $n$ variables \\ \bottomrule
\end{tabular}
\caption{Rendimiento de métodos de diferenciación.}
\end{table}

La diferenciación automática (Autodiff) es el estándar en deep learning. Aplica la regla de la cadena sobre operaciones elementales mediante un grafo computacional.

\subsection{Grafos computacionales}
Un grafo computacional es una representación de las operaciones matemáticas de un modelo como un grafo dirigido acíclico (DAG). Cada nodo representa una operación o variable, y las aristas indican el flujo de datos. Por ejemplo, la expresión $y = (x \cdot w) + b$ se representa como:

\begin{center}
\begin{tikzpicture}[
    node distance=1.5cm,
    io/.style={circle, draw, minimum size=0.6cm},
    op/.style={rectangle, draw, rounded corners, minimum size=0.8cm},
    >=Stealth
]
\node[io] (x) {$x$};
\node[io, below of=x] (w) {$w$};
\node[op, right of=x, xshift=1.5cm] (mul) {$\times$};
\node[io, below of=mul, yshift=0.5cm] (b) {$b$};
\node[op, right of=mul, xshift=1.5cm] (add) {$+$};
\node[io, right of=add, xshift=1.5cm] (y) {$y$};

\draw[->] (x) -- (mul);
\draw[->] (w) -- (mul);
\draw[->] (mul) -- (add);
\draw[->] (b) -- (add);
\draw[->] (add) -- (y);
\end{tikzpicture}
\captionof{figure}{Grafo computacional para $y = x w + b$.}
\end{center}

Para calcular gradientes, se recorre el grafo en sentido inverso (backpropagation) aplicando la regla de la cadena. Los frameworks construyen automáticamente este grafo y proporcionan las derivadas.

\subsection{Anatomía de una función en grafos: ejemplo detallado}
Desarrollemos paso a paso la función $f(x) = \ln(x^2 + 1)$ evaluada en $x=2$, para entender cómo Autodiff calcula la derivada.

\subsubsection{Paso 1: Construcción del grafo computacional}
Identificamos las operaciones intermedias:
\begin{align*}
v_1 &= x^2 \\
v_2 &= v_1 + 1 \\
v_3 &= \ln(v_2) \\
f &= v_3
\end{align*}
El grafo correspondiente es:

\begin{center}
\begin{tikzpicture}[
    node distance=2cm, 
    every node/.style={circle, draw, minimum size=1cm},
    arrow/.style={-Stealth, thick}
]
    \node (x) {$x$};
    \node (v1) [right=of x] {$v_1$};
    \node (v2) [right=of v1] {$v_2$};
    \node (v3) [right=of v2] {$v_3$};

    \draw[arrow] (x) -- node[above, draw=none] {$(\cdot)^2$} (v1);
    \draw[arrow] (v1) -- node[above, draw=none] {$+1$} (v2);
    \draw[arrow] (v2) -- node[above, draw=none] {$\ln(\cdot)$} (v3);
\end{tikzpicture}
\end{center}

\subsubsection{Paso 2: Forward pass (primals)}
Evaluamos los valores numéricos de izquierda a derecha para $x = 2$:
\begin{align*}
v_1 &= x^2 = 4 \\
v_2 &= v_1 + 1 = 5 \\
v_3 &= \ln(5) \approx 1.6094
\end{align*}

\subsubsection{Paso 3: Backward pass (análisis de sensibilidad)}
El objetivo es calcular $\frac{df}{dx}$. Aplicamos la regla de la cadena global:
\[
\frac{df}{dx} = \frac{\partial f}{\partial v_3} \cdot \frac{\partial v_3}{\partial v_2} \cdot \frac{\partial v_2}{\partial v_1} \cdot \frac{\partial v_1}{\partial x}
\]
Calculamos de derecha a izquierda, acumulando el producto de las derivadas locales:

\begin{enumerate}
    \item \textbf{Semilla del gradiente ($\frac{\partial f}{\partial v_3}$):} 
    Iniciamos con la derivada de la salida respecto a sí misma:
    \[
    \frac{\partial f}{\partial v_3} = 1
    \]

    \item \textbf{Eslabón del logaritmo ($\frac{\partial f}{\partial v_2}$):}
    Usamos la derivada local del logaritmo ($1/u$):
    \[
    \frac{\partial f}{\partial v_2} = \frac{\partial f}{\partial v_3} \cdot \frac{\partial v_3}{\partial v_2} = 1 \cdot \frac{1}{v_2} = \frac{1}{5} = 0.2
    \]

    \item \textbf{Eslabón de la suma ($\frac{\partial f}{\partial v_1}$):}
    La derivada local de $(v_1 + 1)$ es $1$:
    \[
    \frac{\partial f}{\partial v_1} = \frac{\partial f}{\partial v_2} \cdot \frac{\partial v_2}{\partial v_1} = 0.2 \cdot 1 = 0.2
    \]

    \item \textbf{Eslabón de la potencia ($\frac{df}{dx}$):}
    La derivada local de $x^2$ es $2x$:
    \[
    \frac{df}{dx} = \frac{\partial f}{\partial v_1} \cdot \frac{\partial v_1}{\partial x} = 0.2 \cdot (2x) = 0.2 \cdot 4 = \mathbf{0.8}
    \]
\end{enumerate}

\subsubsection{Resumen del backward pass}
\begin{itemize}
    \item \textbf{Regla 1:} Siempre empezamos con un $1$ al final del grafo.
    \item \textbf{Regla 2:} En cada flecha hacia atrás, multiplicamos el gradiente que traemos por la derivada local de la operación que cruzamos.
    \item \textbf{Regla 3:} Usamos los valores calculados en el forward pass (primals) para evaluar las derivadas locales.
\end{itemize}

\subsection{La ecuación lineal como grafo tensorial}
En IA, la operación fundamental $y = \mathbf{W}\mathbf{x} + \mathbf{b}$ se visualiza como un grafo donde los parámetros son nodos que permiten el flujo del gradiente. Este es el bloque básico de las capas fully connected.

\begin{center}
\begin{tikzpicture}[
    node distance=1.5cm,
    every node/.style={circle, draw, minimum size=1cm},
    op/.style={rectangle, draw, fill=gray!10},
    arrow/.style={-Stealth}
]
    \node (W) {$\mathbf{W}$};
    \node (x) [below=of W] {$\mathbf{x}$};
    \node (b) [right=of W, xshift=2cm] {$\mathbf{b}$};
    \node[op] (mult) [right=of W, xshift=0.5cm, yshift=-0.7cm] {$\times$};
    \node[op] (add) [right=of mult, xshift=0.5cm, yshift=-0.7cm] {$+$};
    \node (y) [right=of add] {$\mathbf{y}$};

    \draw[arrow] (W) -- (mult);
    \draw[arrow] (x) -- (mult);
    \draw[arrow] (mult) -- (add);
    \draw[arrow] (b) |- (add);
    \draw[arrow] (add) -- (y);
\end{tikzpicture}
\captionof{figure}{Grafo de una capa lineal.}
\end{center}

\subsection{Modos forward y reverse}
\begin{itemize}
    \item \textbf{Modo forward:} Se calculan las derivadas simultáneamente con el valor de la función. Útil cuando hay más salidas que entradas.
    \item \textbf{Modo reverse (backpropagation):} Se calcula primero el valor de la función (forward pass) y luego se propaga el gradiente hacia atrás (reverse pass). Es el más usado en deep learning porque tenemos muchas entradas (parámetros) y una sola salida (pérdida).
\end{itemize}

\section{Implementación manual con NumPy}

Antes de usar frameworks, implementemos descenso de gradiente para un problema simple: regresión lineal con una variable. Usaremos mini-batches y compararemos con la solución analítica.

\subsection{Generación de datos sintéticos}
\begin{lstlisting}[language=Python]
import numpy as np
import matplotlib.pyplot as plt

# Generar datos: y = 2 + 3x + ruido
np.random.seed(42)
X = 2 * np.random.rand(100, 1)
y = 2 + 3 * X + np.random.randn(100, 1)

# Visualizar
plt.scatter(X, y)
plt.xlabel('x')
plt.ylabel('y')
plt.title('Datos sintéticos')
plt.show()
\end{lstlisting}

\subsection{Descenso de gradiente con NumPy}
Implementaremos una clase simple que entrena un modelo lineal con descenso de gradiente mini-batch.

\begin{lstlisting}[language=Python]
class LinearRegressionSGD:
    def __init__(self, lr=0.01, epochs=100, batch_size=32):
        self.lr = lr
        self.epochs = epochs
        self.batch_size = batch_size
        self.w = None
        self.b = None
        self.loss_history = []
    
    def fit(self, X, y):
        n_samples, n_features = X.shape
        # Inicializar parámetros
        self.w = np.random.randn(n_features, 1) * 0.01
        self.b = 0.0
        
        for epoch in range(self.epochs):
            # Mezclar datos
            indices = np.random.permutation(n_samples)
            X_shuffled = X[indices]
            y_shuffled = y[indices]
            
            epoch_loss = 0.0
            for i in range(0, n_samples, self.batch_size):
                # Obtener mini-batch
                X_batch = X_shuffled[i:i+self.batch_size]
                y_batch = y_shuffled[i:i+self.batch_size]
                
                # Forward pass
                y_pred = X_batch @ self.w + self.b
                loss = np.mean((y_pred - y_batch)**2)
                
                # Gradientes
                grad_w = 2 * X_batch.T @ (y_pred - y_batch) / len(X_batch)
                grad_b = 2 * np.mean(y_pred - y_batch)
                
                # Actualizar
                self.w -= self.lr * grad_w.reshape(-1,1)
                self.b -= self.lr * grad_b
                
                epoch_loss += loss * len(X_batch)
            
            epoch_loss /= n_samples
            self.loss_history.append(epoch_loss)
            if epoch % 10 == 0:
                print(f"Epoch {epoch}, loss = {epoch_loss:.4f}")
    
    def predict(self, X):
        return X @ self.w + self.b

# Entrenar
model = LinearRegressionSGD(lr=0.1, epochs=50, batch_size=16)
model.fit(X, y)

# Visualizar pérdida
plt.plot(model.loss_history)
plt.xlabel('Época')
plt.ylabel('Pérdida MSE')
plt.title('Curva de aprendizaje')
plt.show()

# Resultados
print(f"Peso estimado: {model.w[0,0]:.4f}, Intercepto: {model.b:.4f}")
\end{lstlisting}

\section{Introducción a PyTorch}

PyTorch es uno de los frameworks más populares para deep learning. Proporciona tensores con aceleración GPU y diferenciación automática mediante el paquete \texttt{torch.autograd}.

\subsection{Tensores en PyTorch}
\begin{lstlisting}[language=Python]
import torch

# Crear tensores
x = torch.tensor([1.0, 2.0, 3.0])
w = torch.tensor([0.5, -0.2, 0.1], requires_grad=True)
b = torch.tensor(0.0, requires_grad=True)

# Operaciones (se construye el grafo)
y = (x * w).sum() + b
print("y =", y.item())
\end{lstlisting}

\subsection{Cálculo de gradientes con autograd}
\begin{lstlisting}[language=Python]
# Calcular gradientes
y.backward()  # propaga hacia atrás

print("Gradiente de w:", w.grad)
print("Gradiente de b:", b.grad)

# Los gradientes se acumulan; hay que resetearlos si se reutilizan
w.grad.zero_()
b.grad.zero_()
\end{lstlisting}

\subsection{Backpropagation en PyTorch: ejemplo de la función $f(x)=\ln(x^2+1)$}
Podemos verificar el cálculo manual con PyTorch:
\begin{lstlisting}[language=Python]
x = torch.tensor(2.0, requires_grad=True)
y = torch.log(x**2 + 1)
y.backward()
print(f"Gradiente dy/dx en x=2: {x.grad}")  # debe ser 0.8
\end{lstlisting}

\subsection{Entrenamiento de regresión lineal con PyTorch}
Reimplementemos el modelo anterior usando PyTorch.

\begin{lstlisting}[language=Python]
import torch
import torch.nn as nn
import torch.optim as optim

# Convertir datos a tensores de PyTorch
X_t = torch.tensor(X, dtype=torch.float32)
y_t = torch.tensor(y, dtype=torch.float32)

# Definir modelo lineal
model = nn.Linear(1, 1)  # 1 entrada, 1 salida
criterion = nn.MSELoss()
optimizer = optim.SGD(model.parameters(), lr=0.1)

# Entrenamiento
epochs = 50
batch_size = 16
dataset = torch.utils.data.TensorDataset(X_t, y_t)
dataloader = torch.utils.data.DataLoader(dataset, batch_size=batch_size, shuffle=True)

loss_history = []
for epoch in range(epochs):
    epoch_loss = 0.0
    for X_batch, y_batch in dataloader:
        # Forward
        y_pred = model(X_batch)
        loss = criterion(y_pred, y_batch)
        
        # Backward y optimización
        optimizer.zero_grad()
        loss.backward()
        optimizer.step()
        
        epoch_loss += loss.item() * len(X_batch)
    
    epoch_loss /= len(X_t)
    loss_history.append(epoch_loss)
    if epoch % 10 == 0:
        print(f"Epoch {epoch}, loss = {epoch_loss:.4f}")

print(f"Peso: {model.weight.item():.4f}, Intercepto: {model.bias.item():.4f}")
\end{lstlisting}

\section{Introducción a TensorFlow / Keras}

TensorFlow, junto con Keras, ofrece una API de alto nivel para construir y entrenar modelos.

\subsection{Modelo secuencial en Keras}
\begin{lstlisting}[language=Python]
import tensorflow as tf
from tensorflow import keras

# Modelo lineal
model = keras.Sequential([
    keras.layers.Dense(1, input_shape=(1,))
])

model.compile(optimizer=keras.optimizers.SGD(learning_rate=0.1),
              loss='mse')

# Entrenamiento
history = model.fit(X, y, epochs=50, batch_size=16, verbose=0)

print(f"Peso: {model.layers[0].weights[0].numpy()[0,0]:.4f}, Intercepto: {model.layers[0].weights[1].numpy()[0]:.4f}")
\end{lstlisting}

\section{Optimizadores modernos}

Además del SGD básico, existen optimizadores que adaptan la tasa de aprendizaje y utilizan momentos para acelerar la convergencia. Los más comunes son:

\begin{itemize}
    \item \textbf{Momentum}: Acumula una media exponencial de los gradientes para suavizar la actualización.
    \[
    v_{t+1} = \beta v_t + \nabla \mathcal{L}(\mathbf{w}_t), \quad \mathbf{w}_{t+1} = \mathbf{w}_t - \eta v_{t+1}
    \]
    \item \textbf{Adagrad}: Adapta la tasa de aprendizaje por parámetro, usando la suma de cuadrados de gradientes pasados.
    \item \textbf{RMSprop}: Similar a Adagrad pero con media exponencial.
    \item \textbf{Adam (Adaptive Moment Estimation)}: Combina Momentum y RMSprop. Es el más utilizado actualmente.
    \[
    m_t = \beta_1 m_{t-1} + (1-\beta_1) g_t
    \]
    \[
    v_t = \beta_2 v_{t-1} + (1-\beta_2) g_t^2
    \]
    \[
    \hat{m}_t = \frac{m_t}{1-\beta_1^t}, \quad \hat{v}_t = \frac{v_t}{1-\beta_2^t}
    \]
    \[
    \theta_{t+1} = \theta_t - \eta \frac{\hat{m}_t}{\sqrt{\hat{v}_t} + \epsilon}
    \]
    donde $g_t = \nabla_{\theta} \mathcal{L}(\theta_t)$.
\end{itemize}

En PyTorch, podemos usar fácilmente Adam:
\begin{lstlisting}[language=Python]
optimizer = optim.Adam(model.parameters(), lr=0.01)
\end{lstlisting}
En Keras:
\begin{lstlisting}[language=Python]
optimizer = keras.optimizers.Adam(learning_rate=0.01)
\end{lstlisting}

\section{Comparación de implementaciones}

A continuación se muestra una tabla comparativa de las tres implementaciones (NumPy manual, PyTorch, TensorFlow/Keras) para el mismo problema de regresión lineal.

\begin{center}
\begin{tabular}{lccc}
\toprule
\textbf{Característica} & \textbf{NumPy manual} & \textbf{PyTorch} & \textbf{TensorFlow/Keras} \\
\midrule
Cálculo de gradientes & Manual (derivadas) & Automático (autograd) & Automático (GradientTape) \\
Soporte GPU & No & Sí & Sí \\
Facilidad de prototipado & Baja & Alta & Muy alta \\
Control detallado & Total & Total & Medio \\
Curva de aprendizaje & Alta & Media & Baja \\
\bottomrule
\end{tabular}
\end{center}

\section{Notación en papers científicos}

En artículos de deep learning, es común encontrar expresiones como:

\begin{itemize}
    \item Actualización con SGD:
    \[
    \theta_{t+1} = \theta_t - \eta \nabla_{\theta} \mathcal{L}_B(\theta_t)
    \]
    donde $\mathcal{L}_B$ es la pérdida sobre un mini-batch.

    \item Algoritmo de Adam (como se mostró arriba).
    \item Backpropagation en forma vectorial:
    \[
    \frac{\partial \mathcal{L}}{\partial \mathbf{W}^{(l)}} = \delta^{(l)} (\mathbf{a}^{(l-1)})^\top, \quad \delta^{(l)} = (\mathbf{W}^{(l+1)})^\top \delta^{(l+1)} \odot \sigma'(\mathbf{z}^{(l)})
    \]
    \item Diferenciación automática: los papers suelen mencionar que los gradientes se calculan mediante \textit{automatic differentiation} implementada en frameworks como PyTorch o TensorFlow.
\end{itemize}

\section{Ejercicios propuestos}

\begin{enumerate}
    \item Implementa descenso de gradiente estocástico (SGD) puro (tamaño de batch = 1) para regresión lineal y observa la varianza en la curva de pérdida.
    \item Modifica el código de PyTorch para usar el optimizador Adam y compara la velocidad de convergencia con SGD.
    \item Construye el grafo computacional de la función $f(x) = e^{\sin(x^2)}$ y calcula manualmente su derivada en $x=1$ usando el método paso a paso de Autodiff. Luego verifica con PyTorch.
    \item Investiga qué es el ``vanishing gradient'' y cómo los optimizadores modernos ayudan a mitigarlo.
    \item Lee la sección de optimización de un paper de tu interés (por ejemplo, el paper original de Adam) e identifica la notación utilizada para el gradiente y las actualizaciones.
\end{enumerate}

\section{Resumen}

En esta sesión hemos dado el salto de la teoría a la práctica automatizada:
\begin{itemize}
    \item Repasamos las variantes del descenso de gradiente (batch, SGD, mini-batch).
    \item Comparamos los métodos de diferenciación (numérica, simbólica, automática) y profundizamos en Autodiff mediante un ejemplo paso a paso con grafos computacionales.
    \item Implementamos manualmente el descenso de gradiente con NumPy.
    \item Aprendimos los fundamentos de PyTorch (tensores, autograd, entrenamiento) y de TensorFlow/Keras.
    \item Comparamos los optimizadores modernos como Adam.
    \item Vimos cómo se representa esta información en papers científicos.
\end{itemize}
Con estas herramientas, estás listo para abordar problemas más complejos de deep learning, donde los gradientes se calculan automáticamente y los optimizadores se encargan de la actualización de millones de parámetros.

\end{document}